% Generated draft: requires mathematical and source review.
\documentclass[11pt]{article}
\usepackage[margin=25mm]{geometry}
\usepackage{fontspec}
\setmainfont{Latin Modern Roman}
\usepackage{amsmath,amssymb,mathtools,unicode-math}
\setmathfont{Latin Modern Math}
\usepackage{enumitem,needspace,xurl,hyperref}
\setlength{\parindent}{0pt}
\setlength{\parskip}{5pt}
\setlength{\emergencystretch}{3em}
\newcommand{\N}{\mathbb{N}}
\newcommand{\Z}{\mathbb{Z}}
\newcommand{\Q}{\mathbb{Q}}
\newcommand{\R}{\mathbb{R}}
\newcommand{\C}{\mathbb{C}}
\newcommand{\F}{\mathbb{F}}
\newcommand{\A}{\mathbb{A}}
\newcommand{\PP}{\mathbb{P}}
\newcommand{\E}{\mathbb{E}}
\newcommand{\eps}{\varepsilon}
\newcommand{\dd}{\,\mathrm{d}}
\begin{document}
\section*{UnsolvedMath ID: 30004696}
\textbf{Source problem code:} catalogue entry
\subsection{Definitions and mathematical statement}
\textbf{Research question.} What can the use of classical works to legitimize mathematical innovation teach us about the nature of mathematical argument and, more generally, mathematical practice?

This is a historical and interpretive research question. The supplied record does not define a collection of classical works, a historical or geographical scope, or criteria for legitimacy, innovation, mathematical argument, or mathematical practice. These expressions therefore retain their source-language meanings without being assigned formal mathematical definitions. No mathematical domain, hypotheses, quantifiers, boundary conditions, or proposition to prove are specified. A precise theorem or conjecture cannot be extracted without adding unsupported content.

The supplied assessment describes the target as a general account of mathematical argument and practice across historical traditions, framed by the cited workshop material as an open interpretive program. The supplied status \texttt{open} is a curation label, not an independently verified assertion that a particular mathematical proposition remains unproved.
\subsection{Short English statement}
Investigate what appeals to classical mathematical works in legitimizing innovation reveal about mathematical argument and practice. The source poses an interpretive research program; it does not supply enough definitions or scope to formulate a precise mathematical proposition.
\subsection{Sources}
\begin{itemize}
\item Local UnsolvedMath registry: permanent ID \(30004696\), source problem code \texttt{catalogue entry}, title \emph{Classical Authority in Mathematical Innovation and Argument}; collection: \emph{Oberwolfach Reports Open Problems}. The supplied record attributes the question to workshop discussions documented in Oberwolfach Reports and names no individual proposer.
\item Supplied report citation: \emph{Mathematics and its Ancient Classics Worldwide: Translations, Appropriations, Reconstructions, Roles} (2022), Oberwolfach Reports, pp.\ 1347--1406, DOI 10.4171/owr/2021/26. \url{https://doi.org/10.4171/owr/2021/26}. The registry separately supplies proposed year 2021 and the tag \texttt{owr-year:2021}; the relationship between these dates and the citation year has not been independently verified.
\item Supplied literature source: ``Oberwolfach report,'' \url{https://doi.org/10.4171/owr/2021/26}. Literature checked at: \texttt{2026-08-22}. This date and the assessment are supplied dated curation, not independently verified. The supplied assessment labels the question open and describes it as seeking a general account across historical traditions, framed as an open interpretive program. The background explicitly requests independent verification of time-sensitive status and cited claims.
\end{itemize}
\end{document}
